% Work/MAT – nová 19. hodina: další procvičování slovních úloh.
% Návaznost: 17 – dvě množiny; 18 – tři množiny; 20 – opakování celého bloku.
% Zadání vlastní, podle metodiky Realisticky.cz 010306–010307 a učebnice Množiny.tex.
\documentclass[11pt,a4paper]{article}
% Spolecny styl pro ucitelske pripravy MAT (Work/MAT)
% Verze 1.1, 2026-10-23
\NeedsTeXFormat{LaTeX2e}
\ProvidesPackage{priprava}[2026/10/23 v1.1 Sablona ucitelskych priprav MAT]

% Zakladni typografie a cestina
\RequirePackage[T1]{fontenc}
\RequirePackage[utf8]{inputenc}
\RequirePackage[czech]{babel}
\RequirePackage{lmodern}
\RequirePackage{microtype}
\RequirePackage{amsmath,amssymb}
\RequirePackage{xcolor}
\RequirePackage{enumitem}
\RequirePackage[most]{tcolorbox}
\RequirePackage{tabularx,booktabs,array}
\RequirePackage{fancyhdr}
\RequirePackage{needspace}
\RequirePackage{geometry}
\RequirePackage[hidelinks]{hyperref}

\geometry{a4paper,top=17mm,bottom=18mm,left=23mm,right=23mm}
\definecolor{notegray}{gray}{0.39}
\setlength{\parindent}{0pt}
\setlength{\parskip}{2.5pt plus 1pt minus 0.5pt}
\setlist[itemize]{topsep=2pt,itemsep=1pt,parsep=0pt,leftmargin=1.45em}
\setlist[enumerate]{topsep=2pt,itemsep=2pt,parsep=0pt,leftmargin=1.65em}
\setcounter{secnumdepth}{0}
\renewcommand{\arraystretch}{1.12}

% Zahlavi: vlevo rocnik | globalni tema, vpravo cislo: nazev hodiny
% Zapati zustava zcela prazdne.
\pagestyle{fancy}
\fancyhf{}
\renewcommand{\headrulewidth}{0.35pt}
\renewcommand{\footrulewidth}{0pt}
\setlength{\headheight}{14pt}
\newcommand{\nastavPripravu}[4]{%
  \fancyhead[L]{\small #1\enspace\textbar\enspace #2}%
  \fancyhead[R]{\small #3: #4}%
}

% Minimalisticke cernobile boxy. Kazda definice zvlast.
\tcbset{commonlesson/.style={enhanced,breakable=false,colback=white,colframe=black,
  colbacktitle=white,coltitle=black,boxrule=0.45pt,
  arc=0.8mm,outer arc=0.8mm,left=3mm,right=3mm,top=1.6mm,bottom=1.6mm,
  toptitle=1mm,bottomtitle=0.7mm,fonttitle=\bfseries,
  attach boxed title to top left={xshift=2.5mm,yshift=-1.15mm},
  boxed title style={colback=white,colframe=white,boxrule=0pt,left=1mm,right=1mm}}}
\newtcolorbox{zadaniBox}[1][]{commonlesson,title={Zadání},#1}
\newtcolorbox{prikladBox}[1][]{commonlesson,title={Příklad},#1}
\newtcolorbox{definiceBox}[1][]{commonlesson,title={Definice},#1}

% Spolecny box pro zadani a jeho reseni, bez mezery mezi castmi.
% Pouziti:
%   \begin{prikladReseniBox}
%     Zadani ulohy.
%     \tcblower
%     \begin{align*} ... \end{align*}
%   \end{prikladReseniBox}
% Spodni cast ma jemny sedy podklad a delici linku; lze v ni sazet
% vice radku matematiky. Obe casti maji jen jeden spolecny ramecek.
\tcbset{solutionlesson/.style={
  commonlesson,
  enhanced jigsaw,
  breakable,
  colbacklower=black!3,
  segmentation style={solid,draw=black!25,line width=0.35pt},
  fontlower=\small\color{notegray},
  before lower={\noindent\textbf{Řešení:}\enspace},
  leftlower=3mm,rightlower=3mm
}}
\newtcolorbox{prikladReseniBox}[1][]{solutionlesson,title={Příklad},#1}
\newtcolorbox{zadaniReseniBox}[1][]{solutionlesson,title={Zadání},#1}

% Pomucky pro autora pripravy: sedive poznamky; \reseni pro starsi pripravy
\newcommand{\poznamka}[1]{\par\begingroup\small\color{notegray}\itshape
  \textnormal{\textbf{Poznámka pro učitele.}}\enspace #1\par\endgroup}
\newcommand{\reseni}[1]{\par\begingroup\small\color{notegray}
  \textbf{Očekávané řešení:}\enspace #1\par\endgroup}
\newcommand{\blok}[1]{\needspace{5\baselineskip}\subsection*{#1}}
\newcommand{\vyrok}[1]{\uv{#1}}

\usepackage{tikz}
\nastavPripravu{1. ročník}{Množiny}{19}{Slovní úlohy: kombinované příklady}
\begin{document}

\blok{Bez vzorce to jde také: nejdřív přelož text}
Dnes už nezavádíme nový matematický pojem. U každé úlohy si zapiš:
\textbf{(1)} co je celá skupina $U$; \textbf{(2)} co označují
$A,B,C$; \textbf{(3)} co znamená každá známá informace;
\textbf{(4)} jaké počty hledáme. Nezvyklé slovní obraty převáděj na
\emph{disjunktní oblasti} ve Vennově diagramu.

\begin{zadaniReseniBox}[title={D0: rozcvička ve dvojicích},breakable=false]
Vysvětli rozdíly mezi: \uv{oba} a \uv{alespoň jeden};
\uv{právě dva} a \uv{alespoň dva}; \uv{nejvýše dva} a \uv{právě dva}.
\tcblower
U dvou množin \uv{oba} značí průnik; \uv{alespoň jeden} sjednocení.
U tří množin \uv{právě dva} zahrnuje tři oblasti bez středu,
\uv{alespoň dva} zahrnuje také střed. \uv{Nejvýše dva} znamená
vše kromě středu, tedy i \emph{žádnou} nebo právě jednu vlastnost.
\end{zadaniReseniBox}

\blok{Úloha I: zadaný počet v právě jedné množině}
\begin{prikladReseniBox}[title={P1: dvě skupiny, společná část není uvedena}]
Ve skupině $30$ lidí se věnuje šachům $18$ a programování $17$.
\emph{Právě jedné} z aktivit se věnuje $23$ lidí. Kolik se věnuje oběma,
alespoň jedné a žádné?
\tcblower
Je-li $x$ ve společné části, jen šachy $18-x$ a jen programování $17-x$.
\[
(18-x)+(17-x)=23\;\Rightarrow\;35-2x=23\;\Rightarrow\;x=6.
\]
Jen šachy $12$, jen programování $11$; alespoň jedna $12+6+11=29$;
žádná $30-29=1$. Kontrola $12+6=18$, $11+6=17$,
$12+6+11+1=30$.
\end{prikladReseniBox}
\poznamka{Pozor na překlad \uv{právě jedné}: společný průnik do tohoto počtu
nepatří. Nech studenty nejprve nezávisle zapsat výraz pro dvě výlučné oblasti.}

\blok{Úloha II: tři množiny a neznámý střed}
\begin{prikladReseniBox}[title={P2: dopočítání trojitého průniku}]
Ve skupině $50$ zákazníků označují $A,B,C$ užívání tří služeb.
Platí $|A|=27$, $|B|=25$, $|C|=22$,
$|A\cap B|=12$, $|A\cap C|=11$, $|B\cap C|=10$.
\emph{Žádnou} službu nevyužívá $7$ lidí. Kolik využívá všechny tři,
právě dvě, právě jednu? Doplň diagram.
\tcblower
Ve sjednocení je $50-7=43$. Označme $x=|A\cap B\cap C|$.
Z inkluze a exkluze:
\[
43=27+25+22-12-11-10+x=41+x\quad\Rightarrow\quad x=2.
\]
Právě dvě: $AB$ bez $C$: $10$, $AC$ bez $B$: $9$,
$BC$ bez $A$: $8$ (součet $27$).
Pouze $A$: $27-10-9-2=6$;
pouze $B$: $25-10-8-2=5$;
pouze $C$: $22-9-8-2=3$ (součet $14$).
\begin{center}
\begin{tikzpicture}[scale=.74]
\draw[thick] (-3.05,-2.64) rectangle (3.08,2.06);
\draw[thick] (-.85,.38) circle(1.43);\draw[thick] (.85,.38) circle(1.43);
\draw[thick] (0,-.83) circle(1.43);
\node at (-1.76,1.7) {$A$};\node at (1.76,1.7) {$B$};
\node at (0,-2.4) {$C$};\node at (2.85,1.79) {$U$};
\node at (-1.45,.59) {$6$};\node at (1.45,.59) {$5$};
\node at (0,1.14) {$10$};\node at (-.82,-.77) {$9$};
\node at (.82,-.77) {$8$};\node at (0,-.23) {$2$};
\node at (0,-1.63) {$3$};\node at (2.53,-2.07) {$7$};
\end{tikzpicture}
\end{center}
Kontrola: $14+27+2+7=50$;
$6+10+9+2=27$ a obdobně pro $B,C$.
\end{prikladReseniBox}

\blok{Samostatné procvičení bez nápovědy}
\begin{zadaniReseniBox}[title={S1: ověř všechna původní čísla}]
\begin{enumerate}[label=\alph*)]
\item Z $36$ studentů studuje angličtinu $24$, němčinu $18$.
Právě jeden z jazyků studuje $20$. Kolik studuje oba, alespoň jeden a žádný?
\item Ve třídě $32$ žáků sleduje $A$ $19$ žáků, $B$ $17$, $C$ $15$;
$|A\cap B|=8$, $|A\cap C|=7$, $|B\cap C|=6$,
$|A\cap B\cap C|=3$. Kolik nesleduje žádný z nich?
\end{enumerate}
\tcblower
\begin{enumerate}[label=\alph*)]
\item $(24-x)+(18-x)=20$, takže $x=11$.
Jen angličtina $13$, jen němčina $7$, alespoň jeden $31$, žádný $5$.
Kontrola $13+11+7+5=36$.
\item Ve sjednocení $19+17+15-8-7-6+3=33$; údaj je \textbf{nemožný},
protože celek tvoří jen $32$ žáků. Žádný odpovídající diagram neexistuje.
\end{enumerate}
\end{zadaniReseniBox}
\poznamka{Rozvržení 45 min: rozcvička 4 min; P1 10 min; P2 14 min;
samostatná S1 a) 9 min; kontrola a případně S1 b) 6 min; závěr 2 min.
U P2 nech rychlé řešit i diagram, u ostatních stačí přesné mezivýsledky.
S1 b) je diagnostická úloha: výsledek může odhalit nesmyslné zadání.}

\blok{Volitelná rezerva: zkontroluj uskutečnitelnost}
\begin{zadaniReseniBox}[title={R1: co když je počet mimo záporný?},breakable=false]
Ze $20$ lidí se zajímá o $A$ $15$, o $B$ $12$ a obě $4$.
Může existovat popsaná skupina? Jak bys chybu vysvětlil zadavateli?
\tcblower
Ne. Sjednocení by mělo $15+12-4=23$ osob, a to je více než $20$.
Počet mimo sjednocení by vyšel $20-23=-3$, což nelze vyložit jako počet lidí.
\end{zadaniReseniBox}

\blok{Domácí příprava na opakování}
\begin{zadaniReseniBox}[title={D1: dvě a tři množiny}]
\begin{enumerate}[label=\alph*)]
\item Ve skupině $40$ lidí má vlastnost $A$ $22$, $B$ $20$ a právě jednu
z těchto vlastností $24$. Urči průnik a počet bez obou.
\item Pro tři množiny v univerzu $U$ o $40$ prvcích je
$|A|=22$, $|B|=18$, $|C|=16$, $|A\cap B|=8$,
$|A\cap C|=7$, $|B\cap C|=6$, $|A\cap B\cap C|=3$.
Urči počet právě ve dvou a mimo všechny.
\end{enumerate}
\tcblower
\begin{enumerate}[label=\alph*)]
\item $(22-x)+(20-x)=24$, takže $x=9$. Sjednocení $33$, mimo $7$.
\item Právě dvě: $(8-3)+(7-3)+(6-3)=5+4+3=12$.
Sjednocení $22+18+16-8-7-6+3=38$, mimo $40-38=2$.
\end{enumerate}
\end{zadaniReseniBox}
\end{document}
